% problem_id: S001-lemma-characterize-henselian-pair
% source_id: S001 (Stacks Project)
% source_file: more-algebra.tex
% statement_locator: more-algebra.tex lines 2552--2575
% proof_locator: more-algebra.tex lines 2577--2734
% env: lemma
% id_note: label 在本来源内唯一
% pairing: tight (verified)
% status: complete (statement + author proof verbatim + reason + locators)
%
\begin{lemma}
\label{lemma-characterize-henselian-pair}
\begin{reference}
\cite[Chapter XI]{Henselian} and \cite[Proposition 1]{Gabber-henselian}
\end{reference}
Let $(A, I)$ be a pair. The following are equivalent
\begin{enumerate}
\item $(A, I)$ is a henselian pair,
\item given an \'etale ring map $A \to A'$ and an $A$-algebra map
$\sigma : A' \to A/I$, there exists an $A$-algebra map $A' \to A$
lifting $\sigma$,
\item for any finite $A$-algebra $B$ the map $B \to B/IB$ induces
a bijection on idempotents,
\item for any integral $A$-algebra $B$ the map $B \to B/IB$ induces
a bijection on idempotents, and
\item (Gabber) $I$ is contained in the Jacobson radical of $A$ and
every monic polynomial $f(T) \in A[T]$ of the form
$$
f(T) = T^n(T - 1) + a_n T^n + \ldots + a_1 T + a_0
$$
with $a_n, \ldots, a_0 \in I$ and $n \ge 1$ has a root $\alpha \in 1 + I$.
\end{enumerate}
Moreover, in part (5) the root is unique.
\end{lemma}

\begin{proof}
Assume (2) holds. Then $I$ is contained in the Jacobson radical of $A$, since
otherwise there would be a nonunit $f \in A$ congruent to $1$ modulo $I$
and the map $A \to A_f$ would contradict (2). Hence $IB \subset B$
is contained in the Jacobson radical of $B$ for $B$ integral over $A$
because $\Spec(B) \to \Spec(A)$ is closed by
Algebra, Lemmas \ref{algebra-lemma-going-up-closed} and
\ref{algebra-lemma-integral-going-up}.
Thus the map from idempotents of $B$ to idempotents of $B/IB$
is injective by Lemma \ref{lemma-idempotents-determined-modulo-radical}.
On the other hand, since (2) holds, every idempotent
of $B/IB$ lifts to an idempotent of $B$
by Lemma \ref{lemma-lift-idempotent-upstairs}.
In this way we see that (2) implies (4).

\medskip\noindent
The implication (4) $\Rightarrow$ (3) is trivial.

\medskip\noindent
Assume (3). Let $\mathfrak m$ be a maximal ideal and consider the
finite map $A \to B = A/(I \cap \mathfrak m)$. The condition that
$B \to B/IB$ induces a bijection on idempotents implies that
$I \subset \mathfrak m$ (if not, then $B = A/I \times A/\mathfrak m$
and $B/IB = A/I$). Thus we see that $I$ is contained in the Jacobson
radical of $A$. Let $f \in A[T]$ be monic and suppose given a
factorization $\overline{f} = g_0h_0$ with $g_0, h_0 \in A/I[T]$ monic
generating the unit ideal in $A/I[T]$.
Set $B = A[T]/(f)$. Let $\overline{e}$ be the idempotent
of $B/IB$ corresponding to the decomposition
$$
B/IB = A/I[T]/(g_0) \times A/I[T]/(h_0)
$$
of $A$-algebras. Let $e \in B$ be an idempotent lifting $\overline{e}$
which exists as we assumed (3). This gives a product decomposition
$$
B = eB \times (1 - e)B
$$
Note that $B$ is free of rank $\deg(f)$ as an $A$-module.
Hence $eB$ and $(1 - e)B$ are finite locally free $A$-modules.
However, since $eB$ and $(1 - e)B$ have constant rank
$\deg(g_0)$ and $\deg(h_0)$ over $A/I$ we find that the same
is true over $\Spec(A)$. We conclude that
\begin{align*}
f & = \text{CharPol}_A(T : B \to B) \\
& =
\text{CharPol}_A(T : eB \to eB)
\text{CharPol}_A(T : (1 - e)B \to (1 - e)B)
\end{align*}
is a factorization into monic polynomials reducing to the given
factorization modulo $I$. Here $\text{CharPol}_A$ denotes the characteristic
polynomial of an endomorphism of a finite locally free module over $A$.
If the module is free the $\text{CharPol}_A$ is defined as the
characteristic polynomial of the corresponding matrix and in
general one uses Algebra, Lemma \ref{algebra-lemma-standard-covering} to
glue. Details omitted. Thus (3) implies (1).

\medskip\noindent
Assume (1). Let $f$ be as in (5). The factorization of $f \bmod I$
as $T^n$ times $T - 1$ lifts to a factorization $f = gh$ with $g$
and $h$ monic by Definition \ref{definition-henselian-pair}.
Then $h$ has to have degree $1$
and we see that $f$ has a root reducing to $1$ modulo $I$.
Finally, $I$ is contained in the Jacobson radical by
the definition of a henselian pair.
Thus (1) implies (5).

\medskip\noindent
Before we give the proof of the last step, let us show that
the root $\alpha$ in (5), if it exists, is unique. Namely,
due to the explicit shape of $f(T)$, we have
$f'(\alpha) \in 1 + I$ where $f'$ is the derivative of $f$ with
respect to $T$. An elementary argument shows that
$$
f(T) = f(\alpha + T - \alpha) =
f(\alpha) + f'(\alpha) \cdot (T - \alpha) \bmod
(T - \alpha)^2 A[T]
$$
This shows that any other root $\alpha' \in 1 + I$ of $f(T)$
satisfies $0 = f(\alpha') - f(\alpha) = (\alpha' - \alpha)(1 + i)$
for some $i \in I$, so that, since $1 + i$ is a unit in $A$,
we have $\alpha = \alpha'$.

\medskip\noindent
Assume (5). We will show that (2) holds, in other words, that
for every \'etale map $A \to A'$, every section $\sigma : A' \to A/I$
modulo $I$ lifts to a section $A' \to A$.
Since $A \to A'$ is \'etale, the section $\sigma$
determines a decomposition
\begin{equation}
\label{equation-GCHP}
A'/IA' \cong A/I \times C
\end{equation}
of $A/I$-algebras. Namely, the surjective ring map
$A'/IA' \to A/I$ is \'etale by
Algebra, Lemma \ref{algebra-lemma-map-between-etale}
and then we get the desired idempotent by
Algebra, Lemma \ref{algebra-lemma-surjective-flat-finitely-presented}.
We will show that this decomposition lifts to a decomposition 
\begin{equation}
\label{equation-GCHP-want}
A' \cong A'_1 \times A'_2
\end{equation}
of $A$-algebras with $A'_1$ integral over $A$. Then $A \to A'_1$
is integral and \'etale and $A/I \to A'_1/IA'_1$ is an isomorphism,
thus $A \to A'_1$ is an isomorphism by
Lemma \ref{lemma-check-isomorphism-zariski}
(here we also use that an \'etale ring map is flat
and of finite presentation, see Algebra, Lemma \ref{algebra-lemma-etale}).

\medskip\noindent
Let $B'$ be the integral closure of $A$ in $A'$. By
Lemma \ref{lemma-helper-finite-type}
we may decompose 
\begin{equation}
\label{equation-dec-mod-I}
B'/IB' \cong A/I \times C'
\end{equation}
as $A/I$-algebras compatibly with (\ref{equation-GCHP})
and we may find $b \in B'$ that lifts $(1, 0)$ such that
$B'_b \to A'_b$ is an isomorphism. If the decomposition
(\ref{equation-dec-mod-I}) lifts to a decomposition
\begin{equation}
\label{equation-want-2}
B' \cong B'_1 \times B'_2
\end{equation}
of $A$-algebras, then the induced decomposition
$A' = A'_1 \times A'_2$ will give the
desired (\ref{equation-GCHP-want}): indeed, since $b$ is a unit in $B'_1$
(details omitted),
we will have $B'_1 \cong A'_1$, so that $A'_1$ will be integral over $A$. 

\medskip\noindent
Choose a finite $A$-subalgebra $B'' \subset B'$ containing $b$
(observe that any finitely generated $A$-subalgebra of $B'$
is finite over $A$). After enlarging $B''$ we may assume
$b$ maps to an idempotent in $B''/IB''$ producing
\begin{equation}
\label{equation-again-dec-mod-I}
B''/IB'' \cong C''_1 \times C''_2
\end{equation}
Since $B'_b \cong A'_b$ we see that $B'_b$ is of finite type over $A$.
Say $B'_b$ is generated by $b_1/b^n, \ldots, b_t/b^n$ over $A$
and enlarge $B''$ so that $b_1, \ldots, b_t \in B''$. Then
$B''_b \to B'_b$ is surjective as well as injective, hence an isomorphism.
In particular, we see that $C''_1 = A/I$! Therefore $A/I \to C''_1$
is an isomorphism, in particular surjective.
By Lemma \ref{lemma-helper-finite} we can find
an $f(T) \in A[T]$ of the form
$$
f(T) = T^n(T - 1) + a_n T^n + \ldots + a_1 T + a_0
$$
with $a_n, \ldots, a_0 \in I$ and $n \ge 1$ such that $f(b) = 0$.
In particular, we find that $B'$ is a $A[T]/(f)$-algebra.
By (5) we deduce there is a root $a \in 1 + I$ of $f$.
This produces a product decomposition $A[T]/(f) = A[T]/(T - a) \times D$
compatible with the splitting (\ref{equation-dec-mod-I}) of $B'/IB'$.
The induced splitting of $B'$ is then a desired (\ref{equation-want-2}).
\end{proof}
